\def\version{1}
\problemname{The Gill Waltz}
Doris is writing a long email to her family.
The text consists of $K$ different letters, and is $N$ letters long.
Doris does not have a very good memory, so she does not remember where the letters are on the keyboard.
Instead, she uses a variant of the so-called \emph{gill waltz} when typing.

Doris types her text in $K$ rounds, one round for each of the different letters that make up the text.
In the beginning, Doris writes down all occurrences of one of the letters with one of her gills.
After that, she goes back to the beginning of the text and picks a new letter.
Then she writes in all occurrences of this letter, using the right arrow key with her other gill.
After this, she again goes back to the beginning of the text, now to write the third letter, and so on.
She repeats this until she has written the entire text.
This way, she only needs to remember where the key for a single letter is at a time.

Depending on the order she chooses to type the letters in, this can take different amounts of time.
If she is to type the text \texttt{aabbac}, then the order \texttt{a}, \texttt{b}, \texttt{c} requires her to press the right arrow key 7 times.
First, she types \texttt{aaa} without using the right arrow.
Then she goes back to the text, uses the right arrow twice and types \texttt{bb}.
Finally, she goes back to the beginning, presses right five times and types \texttt{c}.

If she instead chose the order \texttt{b}, \texttt{a}, \texttt{c}, she would first type \texttt{bb}.
Then she would use two right presses to type \texttt{aabba}, and finally five more times to type \texttt{aabbac}.

The order \texttt{c}, \texttt{b}, \texttt{a} requires only that she presses the right arrow twice, which is optimal.

\section*{Input}
The first line contains the positive integers $N$ and $K$.
The next line contains a string with $N$ characters, chosen among the $K$ first lowercase letters of the alphabet (\texttt{a}, \texttt{b}, \texttt{c}, \dots).
It is furthermore guaranteed that the string contains exactly $K$ different characters.

\section*{Output}
Output a single number --- the number of times Doris must press the right arrow to type the text.

\section*{Scoring}
Your solution will be tested on a set of test groups. To earn points for a group, you must pass all test cases in that group.

\noindent
\begin{tabular}{| l | l | l | l |}
\hline
Group & Points & Constraints \\ \hline
	1     & 8 & $K = 2$, $N \le 100$ \\ \hline
	2     & 10 & $K \le 7$, $N \le 100$ \\ \hline
	3     & 28 & $K \le 10$, $N \le 1000$ \\ \hline
	4     & 19 & $K \le 10$, $N \le 10^6$ \\ \hline
	5     & 35 & $K \le 18$, $N \le 10^6$ \\ \hline
\end{tabular}
